% Pilot research-direction experiment — moduli of abelian varieties.
% Run with the `planner_critic` workflow.

This is a \emph{pilot} run of the Planner / Worker / Critic loop on a
task that is deliberately off the loop's main specification. The output
is not a proof — it is a research proposal containing \textbf{three
concrete projects} in the area of moduli spaces of abelian varieties,
their compactifications, and intersection theory on them. A panel of
experts in the field will evaluate the result.

Adapt your roles accordingly (see ``Role notes'' below).

\section*{Goal}

Produce a single \LaTeX{} document, \texttt{answer.tex}, which presents
three projects. The intent is to test whether an LLM-driven loop can
identify \textbf{genuinely new and interesting research directions} in
this area, plausibly profiting from the agents' broad cross-field
knowledge to spot opportunities that a domain expert might miss because
they are ``too close'' to the field.

\section*{Hard constraints}

\begin{enumerate}
\item \textbf{Three projects}, each on a \emph{concrete} question or
  potential theorem / calculation / construction. Vague programmatic
  statements (``investigate the behavior of \dots'', ``study
  generalizations of \dots'') are not acceptable.

\item \textbf{Exactly one} of the three projects must target an
  existing open problem in the literature; the other two must propose
  new questions that are themselves motivated and non-obvious. (A ``new
  question'' is not a trivial generalization or specialization of a
  well-known one.)

\item The hoped-for output of each project should plausibly be at the
  level of a paper in \textbf{Compositio Math.}\ or \textbf{Advances in
  Math.} That bar means: a solid theorem with a clear technique, not a
  brief note; not, however, an expected Annals / Inventiones
  blockbuster.

\item \textbf{Length: about 2--3 pages per project; total hard maximum
  11 pages}, including a brief introduction and a bibliography. The
  submitted \texttt{answer.tex} MUST compile to $\le 11$ pages of
  standard \verb|\documentclass[11pt]{article}| output.

\item \textbf{No invented citations.} Every entry in
  \texttt{references.bib} must be a real paper with at least author,
  title, year, and a venue or arXiv id. The Critic has web access for
  spot-checking; do not bluff.
\end{enumerate}

\section*{Required structure}

Each of the three projects must have:

\begin{itemize}
\item \textbf{Central question} --- one to three sentences, sharply
  stated.
\item \textbf{Approach} --- one paragraph: the key idea or ansatz that
  makes this project look feasible. What technique, structural insight,
  or cross-field connection is being leveraged?
\item \textbf{Detailed plan} --- half a page to one page: the concrete
  steps (lemmas to establish, computations to run, intermediate
  reductions). It is acceptable to write ``step $N$ should go through
  by an argument analogous to $X$ in [Cite]'' when $X$ is genuinely
  analogous; it is \emph{not} acceptable to wave at unspecified
  machinery.
\item \textbf{Why this is plausibly tractable} --- the strongest reason
  a panel member should believe the project has a real chance.
\item \textbf{Risks / open questions} --- what could derail it, and
  what the fallback partial result would still look like.
\end{itemize}

The document should open with a short ($\le 1$ page) framing section
explaining the perspective taken across the three projects (e.g.\ ``these
three projects share a focus on \dots''). It should close with a unified
\texttt{references.bib} bibliography.

\section*{Evaluation criteria (the panel will grade on)}

\begin{enumerate}
\item \textbf{Originality} --- would an active researcher in moduli of
  abelian varieties / compactifications / intersection theory
  recognize each project as genuinely new and interesting, or as
  something already in the air / well-known?

\item \textbf{Feasibility} --- does the proposed approach plausibly
  close, or at least produce a clear partial result, given the
  techniques cited?

\item \textbf{Factual correctness} --- every statement of an existing
  theorem, citation, or implication used in the plan must be correct.
\end{enumerate}

\section*{Role notes (adaptation from the default proof-correctness loop)}

This is a research-proposal task, not a proof task. The default loop
expects the Worker to write a proof and the Critic to find errors in
it. Here:

\textbf{Planner.} You take the lead. Curate a \emph{portfolio} of
candidate projects across rounds --- keep more candidates in
\texttt{research\_notes.tex} than you ultimately ship; rank and prune
across rounds. Cast a wide net using web search: look beyond the
immediate subfield (Hodge theory, motivic / motivic homotopy, derived
categories, $p$-adic Hodge theory, condensed methods, tropical /
non-archimedean geometry, formal-group / arithmetic-geometric tools,
representation theory of automorphic forms touching $\mathcal{A}_g$,
\dots) for transferable techniques that domain insiders may have
overlooked. By the final rounds, narrow to the strongest three.

\textbf{Worker.} Your job is \emph{not} to write proofs. Per round, do
one or more of: literature dives (download / skim arXiv papers,
populate \texttt{references.bib}, summarise key results in
\texttt{research\_notes.tex}); sanity-check toy computations (Sage /
Python / SymPy --- e.g.\ verify a dimension count, a low-genus
intersection number, a small Hodge integral); first-pass feasibility
analysis (chase implications of the cited results to see whether the
proposed approach hits an obvious wall). Update \texttt{answer.tex}
paragraph by paragraph. The four canonical files (\texttt{answer.tex},
\texttt{research\_notes.tex}, \texttt{references.bib},
\texttt{report.md}) form the contract; respect it.

\textbf{Critic.} Reinterpret the BLOCKER / CONCERN / NIT taxonomy for
this task as follows:

\begin{itemize}
\item \textbf{BLOCKER} = a plainly false mathematical claim, an
  invented or miscited reference, a project description whose stated
  approach contradicts something in the literature, or a project that
  the Critic believes (with a citation) is already a published result.
\item \textbf{CONCERN} = an unsupported feasibility argument, a
  hand-wave at a load-bearing step, a vague approach that hasn't been
  concretised, an originality concern (the question seems to be already
  widely considered in the area), an over-broad or multi-paper-sized
  scope packaged as a single project.
\item \textbf{NIT} = style, organisation, redundancy, things that don't
  affect the panel's grade.
\end{itemize}

Use web search to spot-check the citations and to verify that proposed
open questions are not already known.

The early-stop ``answer ready'' gate likely will not fire on this kind
of task; expect to use all rounds. That is fine.

\section*{Tools available to the Worker}

\begin{itemize}
\item Python 3 (\texttt{numpy}, \texttt{scipy}, \texttt{sympy},
  \texttt{matplotlib}, \texttt{networkx}, \texttt{pandas});
\item Internet access (\texttt{curl} / \texttt{wget} for arXiv, web
  search via the host shell);
\item \texttt{pdflatex}, \texttt{bibtex}, \texttt{biber};
\item The full workspace persists across rounds.
\end{itemize}

\section*{Pilot framing --- please be self-aware}

This is the first time we are running this loop on a research-direction
task. The default prompts may push you toward proof-checking habits.
When the situation calls for it, prefer mathematical \emph{judgement}
(likelihood the approach closes, novelty of the question) over formal
verification (which is mostly not applicable here). When you make a
judgement call, mark it explicitly in \texttt{research\_notes.tex} so
the panel can see your reasoning. Where you are confident vs where you
are speculating should be clear from the document.
